\documentclass[11pt,a4paper]{moderncv}

\moderncvstyle{classic}
\moderncvcolor{blue}

\usepackage[utf8]{inputenc}
\usepackage[T1]{fontenc}
\usepackage[spanish]{babel}
\usepackage[scale=0.9]{geometry}
\usepackage{graphicx}
\usepackage{fontawesome5}
\renewcommand*{\addresssymbol}{\faIcon{map-marker}\hspace{0.5em}}
\renewcommand*{\emailsymbol}{\faIcon{envelope}\hspace{0.5em}}
\renewcommand*{\homepagesymbol}{\faIcon{linkedin}\hspace{0.5em}}

\name{Miguel Ángel}{Marcial Martínez}
\title{Ingeniero en Computación\\
\small Maestrante en Ciencias de Datos}

\address{Oaxaca, México}{}{}
\email{fitorec@mundosica.com}
\homepage{linkedin.com/in/fitorec}
\extrainfo{\faIcon{phone} +52 951 115 9304}
\photo[64pt][0.4pt]{img/perfil}

\makeatletter
\renewcommand*{\namefont}{\fontsize{20}{5}\mdseries\upshape}
\makeatother

\begin{document}
\makecvtitle

% ===============================
% Perfil profesional
% ===============================

\section{Perfil Profesional}

Ingeniero en Computación, egresado de la Universidad Tecnológica de la Mixteca (UTM), con experiencia en desarrollo de software, infraestructura de datos y divulgación tecnológica. Mi trabajo se ha centrado en el análisis y la construcción de sistemas digitales que impactan procesos sociales, institucionales y de toma de decisiones, con especial interés en la transparencia, la gobernanza de datos y el uso responsable de la tecnología.



He participado como docente y ponente en diversas instituciones de educación superior, impartiendo cursos, talleres y conferencias relacionados con ingeniería, programación, ciencia de datos y reflexión crítica sobre la tecnología. Mi enfoque combina la práctica técnica con una visión interdisciplinaria que integra dimensiones sociales, éticas y filosóficas del desarrollo tecnológico.


Actualmente trabajo en la optimización de infraestructura crítica de datos en México, así como en el diseño y la liberación pública de algoritmos de identidad digital, promoviendo principios de software libre, auditabilidad y soberanía tecnológica. Considero que los sistemas que sostienen funciones esenciales de la vida pública deben ser abiertos, comprensibles y verificables.


Mis intereses incluyen la infraestructura digital pública, la identidad digital, la gobernanza de datos, el software de código abierto y la relación entre tecnología, poder y sociedad. Busco contribuir al desarrollo de tecnologías que no solo sean eficientes, sino también socialmente responsables y orientadas al bien común.


% ===============================
% Secciones
% ===============================
\section{Experiencia Profesional}

\cventry
{2022--2024}
{Desarrollador Principal / Seguridad de Aplicaciones}
{MRCI -- Plataforma Cyborg (Entidad financiera)}
{Remoto / México}
{}
{
\begin{itemize}
  \item Desarrollo principal de la plataforma financiera \textit{Cyborg}.
  \item Implementación de algoritmos de riesgo comercial y prevención de lavado de dinero (PLD).
  \item Revisión y aplicación de políticas de seguridad bajo lineamientos OWASP.
  \item Automatización de despliegues mediante GitLab CI/CD en infraestructura Linode.
\end{itemize}
}

\cventry
{2015--Actualidad}
{Socio, Project Manager y Responsable Técnico}
{HagamosWeb}
{México}
{}
{
\begin{itemize}
  \item Gestión integral de proyectos de desarrollo de software.
  \item Definición y priorización de requerimientos junto con Product Owner.
  \item Implementación de metodologías ágiles (Scrum).
  \item Liderazgo técnico y coordinación de equipos multidisciplinarios.
\end{itemize}
}

\cventry
{2013}
{Desarrollador de Sistemas}
{Carnes Cruvi}
{México}
{}
{
\begin{itemize}
  \item Desarrollo de sistema de gestión de clientes y precios diferenciados.
  \item Implementación de plataforma web para restaurantes y carnicerías.
\end{itemize}
}

\cventry
{2010--Actualidad}
{Desarrollador de Sistemas}
{Multiversidad Mundo Sin Fronteras}
{México}
{}
{
\begin{itemize}
  \item Desarrollo del sistema de pagos, recargos y condonaciones de colegiaturas.
  \item Implementación de plataforma académica en línea.
\end{itemize}
}

\cventry
{2010}
{Socio y Desarrollador Principal}
{Mundosica.com}
{México}
{}
{
\begin{itemize}
  \item Desarrollo de productos web y soluciones tecnológicas.
  \item Participación en emprendimiento tecnológico desde etapas tempranas.
\end{itemize}
}

\cventry
{2008--2010}
{Programador y Analista de Seguridad}
{Universidad Tecnológica de la Mixteca (UTM)}
{Oaxaca, México}
{}
{
\begin{itemize}
  \item Desarrollo del sistema SIDEP para el área de posgrado.
  \item Configuración y administración de servidores Linux (Gentoo).
  \item Análisis de seguridad en aplicaciones web y redes institucionales.
  \item Implementación de prácticas de seguridad basadas en OWASP.
\end{itemize}
}

\cventry
{2023--Actualidad}
{Desarrollador -- Proyecto Personal}
{MEXA}
{Open Source}
{}
{
\begin{itemize}
  \item Análisis e implementación de algoritmos oficiales de tramitología mexicana (CURP).
  \item Desarrollo de pruebas automatizadas con datasets de más de 15,000 usuarios.
  \item Despliegue continuo usando GitHub Actions y GitHub Pages.
  \item Proyecto disponible en: \url{https://fitorec.github.io/mexa-py/}
\end{itemize}
}

\section{Divulgación Ponencias y Talleres}

\cventry
{2023}
{Ponente}
{Instituto Tecnológico del Valle de Oaxaca (ITVO)}
{Oaxaca, México}
{}
{
\begin{itemize}
  \item Ponencia: \textit{``Chamilo como sistema de gestión de información en la Universidad Multiversidad Mundo Sin Fronteras''}.
  \item Exposición de experiencias prácticas en la implementación de plataformas educativas.
\end{itemize}
}

\cventry
{2021}
{Instructor}
{Instituto Tecnológico del Valle de Oaxaca (ITVO)}
{Oaxaca, México}
{}
{
\begin{itemize}
  \item Taller de Arduino y microcontroladores ESP8266.
  \item Introducción a sistemas embebidos con conectividad WiFi.
  \item Aplicación práctica en proyectos reales (invernadero automatizado con Raspberry Pi).
\end{itemize}
}

\cventry
{2019}
{Instructor}
{Instituto Tecnológico del Valle de Oaxaca (ITVO)}
{Oaxaca, México}
{}
{
\begin{itemize}
  \item Taller de desarrollo web con jQuery.
  \item Enfoque práctico en interacción cliente-servidor y manipulación del DOM.
  \item Material del curso disponible en línea.
\end{itemize}
}

\cventry
{2017}
{Ponente}
{Universidad Tecnológica de la Mixteca (UTM)}
{Huajuapan de León, Oaxaca}
{}
{
\begin{itemize}
  \item Simposium de Software IX.
  \item Ponencia: \textit{Experiencia y visión de egresados}.
  \item Presentación de una metodología MVC aplicada a aplicaciones web modernas.
\end{itemize}
}

\cventry
{2012}
{Ponente}
{Universidad José Vasconcelos}
{Oaxaca, México}
{}
{
\begin{itemize}
  \item Ponencia: \textit{``Software Libre: Rebeldia y Libertad''} en el marco del FLISOL Oaxaca 2012.
  \item Difusión de los principios filosóficos y técnicos del software libre.
  \item Repositorio: \url{https://github.com/fitorec/softwareLibre_rebeldia_y_libertad}
\end{itemize}
}

\cventry
{2012}
{Instructor}
{Universidad Tecnológica de la Mixteca (UTM)}
{Huajuapan de León, Oaxaca}
{}
{
\begin{itemize}
  \item Simposium de Software VIII.
  \item Taller: \textit{Desarrollo Web MVC con Backbone.js}.
  \item Presentación de una metodología MVC aplicada a aplicaciones web modernas.
\end{itemize}
}


\cventry
{2008--Actualidad}
{Miembro y Divulgador}
{Comunidad de Software Libre en Oaxaca}
{Oaxaca, México}
{}
{
\begin{itemize}
  \item Cofundador y participante activo en la comunidad de software libre en Oaxaca.
  \item Organización y participación en eventos, talleres y encuentros tecnológicos.
  \item Promoción de la cultura del software libre, el conocimiento abierto y la colaboración.
\end{itemize}
}
\section{Formación Académica}

\cventry
{2024}
{Licenciatura en Ingeniería en Computación}
{CENEVAL -- Acuerdo 286 (Acreditación por Experiencia Laboral)}
{México}
{}
{
\begin{itemize}
  \item Titulación obtenida mediante el Acuerdo 286 del CENEVAL.
  \item Reconocimiento de conocimientos y competencias adquiridas por experiencia profesional.
  \item Cédula profesional: \textbf{14861759}.
\end{itemize}
}

\cventry
{--}
{Ingeniería en Computación}
{Universidad Tecnológica de la Mixteca (UTM)}
{Huajuapan de León, Oaxaca}
{}
{
\begin{itemize}
  \item 100\% de créditos concluidos.
  \item Formación sólida en computación, desarrollo de software y fundamentos de ingeniería.
\end{itemize}
}

\cventry
{En curso}
{Maestría en Ciencia de Datos}
{Programa de Posgrado}
{México}
{}
{
\begin{itemize}
  \item Formación avanzada en análisis de datos, estadística, aprendizaje automático e inteligencia artificial.
  \item Enfoque en aplicaciones prácticas y p

\section{Habilidades Técnicas}

\cvitem{Lenguajes de programación}{
Python, JavaScript, PHP, Bash, SQL, C/C++, Java, Ruby, Perl
}

\cvitem{Desarrollo Web}{
HTML5, CSS3, Sass, Less, JavaScript, jQuery, Backbone.js, Node.js,
Flask, WordPress
}

\cvitem{Bases de datos}{
MySQL, PostgreSQL, MongoDB, PL/SQL
}

\cvitem{Ciencia de Datos e IA}{
Análisis de datos, limpieza y preparación de datos,
algoritmos de clasificación y regresión,
fundamentos de aprendizaje automático,
visualización de datos
}

\cvitem{Sistemas y DevOps}{
Linux (Debian, Gentoo), Unix, MacOS, Windows,
Git, GitLab CI/CD, GitHub Actions,
automatización de despliegues
}

\cvitem{Seguridad}{
Buenas prácticas OWASP,
análisis de vulnerabilidades,
seguridad en aplicaciones web,
prevención de riesgos en sistemas financieros
}

\cvitem{Herramientas}{
Docker (nivel intermedio),
WPS Office, LibreCAD, Geany,
entornos de desarrollo multiplataforma
}

\cvitem{Metodologías}{
Scrum, gestión de proyectos,
trabajo colaborativo,
documentación técnica
}

\cvitem{Hardware y Sistemas Embebidos}{
Arduino, ESP8266,
Raspberry Pi,
integración hardware--software
}

\section{Habilidades Blandas}

\cvitem{}{
Comunicación efectiva,
liderazgo técnico,
pensamiento analítico,
resolución de problemas,
gestión del tiempo,
aprendizaje autónomo,
trabajo en equipo
}

\section{Contribuciones y Proyectos}

\cventry
{2025--Actualidad}
{Colaborador y Desarrollador}
{Identidad -- Proyecto MxOS}
{Open Source}
{}
{
\begin{itemize}
  \item Desarrollo de bibliotecas en \textbf{Rust} para la validación y generación de identificadores oficiales mexicanos:
  CURP, NSS y RFC.
  \item Diseño de soluciones eficientes, auditables y seguras para procesos de identidad digital.
  \item Enfoque en estandarización, software libre y soberanía tecnológica.
  \item Repositorio: \url{https://github.com/Mx-OS/identidad}
\end{itemize}
}

\cventry
{2025--Actualidad}
{Colaborador}
{MX-OS -- Sistema Operativo Linux Mexicano}
{Open Source}
{}
{
\begin{itemize}
  \item Participación en la iniciativa estratégica para el desarrollo de un sistema operativo Linux mexicano.
  \item Diseño del arte gráfico, logotipo y manual de identidad visual del proyecto.
  \item Colaboración activa con la comunidad y definición de lineamientos visuales.
\end{itemize}
}

\cventry
{2023--Actualidad}
{Autor y Desarrollador Principal}
{MEXA}
{Proyecto Personal / Open Source}
{}
{
\begin{itemize}
  \item Desarrollo de un validador y generador de campos para tramitología mexicana
  (CURP, RFC, NSS, CLABE, entre otros).
  \item Implementación de pruebas automatizadas y despliegue continuo.
  \item Proyecto disponible en:
  \url{https://github.com/fitorec/mexa-py} \\
  Página del proyecto: \url{https://fitorec.github.io/mexa-py/}
\end{itemize}
}

\cventry
{2018}
{Colaborador}
{Chamilo LMS}
{Open Source}
{}
{
\begin{itemize}
  \item Desarrollo de nuevas funcionalidades para la gestión de cursos, usuarios y herramientas de aprendizaje.
  \item Refactorización masiva del código, modernizando la arquitectura y adoptando mejores prácticas en PHP.
  \item Optimización de consultas a base de datos y mejora del rendimiento general del sistema.
  \item Corrección de bugs críticos y vulnerabilidades de seguridad.
  \item Aseguramiento de compatibilidad con versiones recientes de PHP y dependencias.
  \item Repositorio: \url{https://github.com/fitorec/chamilo-lms/commits?author=fitorec}
\end{itemize}
}

\cventry
{2015}
{Colaborador}
{Javalite}
{Open Source}
{}
{
\begin{itemize}
  \item Actualización de dependencias y herramientas de construcción (Maven, JUnit).
  \item Correcciones de mantenimiento en configuración, documentación y scripts de build.
  \item Ajustes para mejorar compatibilidad y estabilidad con nuevas versiones de librerías.
  \item Repositorio: \url{https://github.com/javalite/javalite/commits?author=fitorec}
\end{itemize}
}

\cventry
{2015}
{Colaborador}
{Joomla! CMS}
{Open Source}
{}
{
\begin{itemize}
  \item Corrección de bugs en componentes del core relacionados con sesiones, autenticación y validación de datos.
  \item Implementación de mejoras de seguridad y estabilidad en módulos críticos.
  \item Ajustes de compatibilidad con nuevas versiones de PHP y librerías externas.
  \item Refactorizaciones menores para mantener estándares de codificación.
\end{itemize}
}

\cventry
{2012}
{Colaborador}
{LibreCAD}
{Open Source}
{}
{
\begin{itemize}
  \item Contribuciones en C++ para mejorar el proceso de instalación en sistemas Linux.
  \item Ajustes y mejoras en la traducción al español mexicano.
  \item Repositorio: \url{https://github.com/LibreCAD/LibreCAD/commits?author=fitorec}
\end{itemize}
}

\cventry
{2011--2012}
{Colaborador}
{CakePHP Framework}
{Open Source}
{}
{
\begin{itemize}
  \item Correcciones y mejoras en la documentación y comentarios del código.
  \item Solución de bugs relacionados con validaciones y comportamiento de componentes.
  \item Refactorizaciones menores y ajustes en pruebas para mejorar la estabilidad del framework.
  \item Repositorio: \url{https://github.com/cakephp/cakephp/commits?author=fitorec}
\end{itemize}
}


\cventry
{2010--2022}
{Autor y Desarrollador Principal}
{DiaSQL}
{Proyecto Personal / Open Source}
{}
{
\begin{itemize}
  \item Plugins para el modelador de diagramas de base de datos.
  \item Permite exportar tus diagramas Dia a esquemas SQL.
  \item Proyecto disponible en:
  \url{https://github.com/fitorec/diasql}
\end{itemize}
}

\section{Idiomas}

\cvitemwithcomment{Español}{Nativo}{Lengua materna}

\cvitemwithcomment{Inglés}{Intermedio}{
Lectura y escritura técnica,
comprensión de documentación,
comunicación básica oral
}



\end{document}
